%-------------------------
% Resume - Sukhrob Ilyosbekov
% PhD / Research-focused CV
%-------------------------
\documentclass[letterpaper,11pt]{article}

\usepackage[utf8]{inputenc}
\usepackage[T1]{fontenc}
\usepackage{charter}
\usepackage[margin=0.5in]{geometry}
\usepackage{titlesec}
\usepackage{enumitem}
\usepackage{hyperref}
\usepackage{fancyhdr}
\usepackage{xcolor}
\usepackage{fontawesome5}
\usepackage{tabularx}

% Colors
\definecolor{accent}{HTML}{2B5797}
\definecolor{darktext}{HTML}{1A1A1A}
\definecolor{lighttext}{HTML}{444444}
\definecolor{linkcolor}{HTML}{0563C1}
\definecolor{rulecolor}{HTML}{2B5797}

% Hyperlink setup
\hypersetup{
    colorlinks=true,
    linkcolor=linkcolor,
    urlcolor=linkcolor,
    pdfborderstyle={/S/U/W 0},
}

% Page style
\pagestyle{fancy}
\fancyhf{}
\renewcommand{\headrulewidth}{0pt}
\renewcommand{\footrulewidth}{0pt}

% Section formatting - colored rule
\titleformat{\section}{
    \color{darktext}\large\bfseries\scshape
}{}{0em}{}[{\color{rulecolor}\vspace{2pt}\titlerule\vspace{1pt}}]
\titlespacing*{\section}{0pt}{7pt}{4pt}

% Tighter list spacing
\setlist[itemize]{leftmargin=1.4em, itemsep=0pt, parsep=0pt, topsep=1pt}
\setlist[description]{leftmargin=0pt, labelsep=0.5em, font=\bfseries\color{darktext}, style=unboxed, itemsep=0pt}

% --- Education: university + date on first line, details as bullets ---
\newcommand{\educationEntry}[2]{%
    \vspace{1pt}\noindent
    \begin{tabular*}{\textwidth}{@{}l@{\extracolsep{\fill}}r@{}}
        \textbf{\color{accent}#1} & \textit{\color{lighttext}#2} \\
    \end{tabular*}\vspace{-3pt}
}

% --- Experience: role | company | location, date right-aligned ---
\newcommand{\experienceEntry}[4]{%
    \vspace{3pt}\noindent
    \begin{tabular*}{\textwidth}{@{}l@{\extracolsep{\fill}}r@{}}
        \textbf{\color{darktext}#1} {\color{lighttext}\textbar} {\color{accent}#2} {\color{lighttext}\textbar} {\color{lighttext}#3} & \textit{\color{lighttext}#4} \\
    \end{tabular*}\vspace{-4pt}
}

% --- Project: name + GitHub link, stack on second line ---
\newcommand{\projectEntry}[3]{%
    \vspace{2pt}\noindent
    \begin{tabular*}{\textwidth}{@{}l@{\extracolsep{\fill}}r@{}}
        {\textbf{\color{accent}#1}} & {\small\href{#3}{\color{linkcolor}GitHub}} \\
    \end{tabular*}\vspace{-3pt}\\
    {\small\color{lighttext}\textit{#2}}\vspace{-2pt}
}

% --- Project: GitHub + Paper links, stack on second line ---
\newcommand{\projectEntryPaper}[4]{%
    \vspace{2pt}\noindent
    \begin{tabular*}{\textwidth}{@{}l@{\extracolsep{\fill}}r@{}}
        {\textbf{\color{accent}#1}} & {\small\href{#3}{\color{linkcolor}GitHub} \enspace\color{lighttext}\textbar\enspace \href{#4}{\color{linkcolor}Paper}} \\
    \end{tabular*}\vspace{-3pt}\\
    {\small\color{lighttext}\textit{#2}}\vspace{-2pt}
}

% --- Project: Paper link only, stack on second line ---
\newcommand{\projectEntryPaperOnly}[3]{%
    \vspace{2pt}\noindent
    \begin{tabular*}{\textwidth}{@{}l@{\extracolsep{\fill}}r@{}}
        {\textbf{\color{accent}#1}} & {\small\href{#3}{\color{linkcolor}Paper}} \\
    \end{tabular*}\vspace{-3pt}\\
    {\small\color{lighttext}\textit{#2}}\vspace{-2pt}
}

\begin{document}

%----------HEADING----------
\begin{center}
    {\Huge\bfseries\color{darktext} Sukhrob Ilyosbekov}\vspace{4pt}\\
    {\large\color{accent} Computer Vision \enspace\textperiodcentered\enspace Deep Learning \enspace\textperiodcentered\enspace Explainable AI}\vspace{8pt}\\
    {\color{lighttext}
    \faMapMarker*\enspace Portland, ME \quad
    \faPhone\enspace(857) 867-1942 \quad
    \faEnvelope\enspace\href{mailto:ilyosbekov.s@northeastern.edu}{ilyosbekov.s@northeastern.edu}
    }\vspace{3pt}\\
    {\color{lighttext}
    \faGlobe\enspace\href{https://suxrobgm.net/research}{suxrobgm.net/research} \quad
    \faGraduationCap\enspace\href{https://scholar.google.com/citations?user=p7ujRHoAAAAJ}{Google Scholar} \quad
    \faLinkedin\enspace\href{https://www.linkedin.com/in/suxrobgm}{linkedin.com/in/suxrobgm} \quad
    \faGithub\enspace\href{https://github.com/suxrobgm}{github.com/suxrobgm}
    }
\end{center}

%----------RESEARCH INTERESTS----------
\section{Research Interests}
I'm drawn to the gap between a vision model that scores well and one people can actually rely on. I first ran into it in medical imaging, where a skin-lesion classifier can be accurate and still be looking at the wrong part of the image, with no easy way for a clinician to notice. The same question has followed me into cell microscopy and generative image editing. In a PhD I'd like to work on representation learning and evaluation for medical and scientific images, where labels are expensive to get and a confident mistake matters more than a leaderboard number. Building models for patient care in industry is a big part of why I care about this.

%----------EDUCATION----------
\section{Education}

\educationEntry{Northeastern University, Boston, MA}{January 2024 -- May 2026}
\begin{itemize}
    \item Master of Science in Computer Science; GPA 4.0/4.0
    \item Coursework: Advanced Perception, Computer Vision, Machine Learning
\end{itemize}

\educationEntry{Suffolk University, Boston, MA}{September 2021 -- May 2023}
\begin{itemize}
    \item Bachelor of Science in Computer Science; Cum Laude; GPA 3.5/4.0
    \item Transferred from INTI International College, Malaysia (2019 -- 2021) and Tashkent University of Information Technologies, Uzbekistan (2017 -- 2019)
\end{itemize}

%----------PUBLICATIONS----------
\section{Publications}
\begin{enumerate}[leftmargin=1.4em, itemsep=2pt, label={[\arabic*]}]
    \item \textbf{S.\ Ilyosbekov}, S.\ Gajjar, and R.\ Jin, ``MorphoCLIP: Text-Supervised Contrastive Learning for Perturbation Matching in Cell Painting Images,'' \textit{arXiv preprint}, arXiv:2608.22690, 2026. \href{https://arxiv.org/abs/2608.22690}{\color{linkcolor}[Paper]}
    % ON SUBMISSION: swap for the line below. Targets: LMRL @ NeurIPS, ML4H, MIDL.
    % \item \textbf{S.\ Ilyosbekov}, S.\ Gajjar, and R.\ Jin, ``MorphoCLIP: Text-Supervised Contrastive Learning for Perturbation Matching in Cell Painting Images,'' \textit{under review}, VENUE 2027. \href{https://arxiv.org/abs/2608.22690}{\color{linkcolor}[Paper]}
    \item \textbf{S.\ Ilyosbekov}, ``Localize, Don't Beautify: Client-Side Control of Image-Editing APIs for Cosmetic Surgery Previews,'' \textit{arXiv preprint}, arXiv:2608.02841, 2026. \href{https://arxiv.org/abs/2608.02841}{\color{linkcolor}[Paper]}
    % ON SUBMISSION: targets: CVPR/ICCV generative-eval workshop, WACV.
    % \item \textbf{S.\ Ilyosbekov}, ``Localize, Don't Beautify: Client-Side Control of Image-Editing APIs for Cosmetic Surgery Previews,'' \textit{under review}, VENUE 2027. \href{https://arxiv.org/abs/2608.02841}{\color{linkcolor}[Paper]}
    \item \textbf{S.\ Ilyosbekov}, ``MelanomaNet: Explainable Deep Learning for Multi-Class Skin Lesion Classification,'' \textit{arXiv preprint}, arXiv:2512.09289, 2025. \href{https://arxiv.org/abs/2512.09289}{\color{linkcolor}[Paper]}
    % ON SUBMISSION: targets: ISBI, MIDL, ISIC @ MICCAI.
    % \item \textbf{S.\ Ilyosbekov}, ``MelanomaNet: Explainable Deep Learning for Multi-Class Skin Lesion Classification,'' \textit{under review}, VENUE 2027. \href{https://arxiv.org/abs/2512.09289}{\color{linkcolor}[Paper]}
\end{enumerate}

%----------AWARDS & HONORS----------
\section{Awards \& Honors}
% TODO: add the year to each award so the list is consistent.
\begin{itemize}[itemsep=2pt]
    \item \textbf{Associate Member,} Sigma Xi, The Scientific Research Honor Society, 2026
    \item \textbf{Outstanding Project Award,} Bookshelf Scanner, Northeastern University
    \item \textbf{National Physics Olympiad, 1st Place,} Uzbekistan; first student from my high school to win nationally
    \item \textbf{University Coding Competition Winner,} Tashkent University of Information Technologies
\end{itemize}

%----------RESEARCH----------
\section{Research}

\projectEntryPaper{MorphoCLIP: Text-Supervised Cell Painting Retrieval}{PyTorch, DINOv3, BioClinical ModernBERT, contrastive learning, CPJUMP1}
{https://github.com/suxrobgm/morphoclip}
{https://arxiv.org/abs/2608.22690}
\begin{itemize}
    \item Learns a shared space for Cell Painting images and text, so a compound or gene perturbation can be retrieved from a written description of it, and the other way round.
    \item DINOv3 and BioClinical ModernBERT stay frozen and only the projection heads train, which fits on a single consumer GPU.
    \item Evaluated on CPJUMP1: 51 plates, 3M+ images, 303 compounds, and 160 genes. Batch correction in embedding space keeps the model from simply recognizing plates.
    \item First author, with S.\ Gajjar and R.\ Jin. I led model design, training, and evaluation.
\end{itemize}

\projectEntryPaperOnly{Localize, Don't Beautify: Client-Side Control of Image-Editing APIs}{Image-editing APIs, inpainting models, ArcFace, prompt engineering}
{https://arxiv.org/abs/2608.02841}
\begin{itemize}
    \item Studies how far a client can control a black-box editing API, using cosmetic-surgery previews as the test case.
    \item I ran prompt steering, masked compositing, and model-based inpainting through six commercial editors on 196 facelift and rhinoplasty edits.
    \item Masked compositing, the simplest of the three, kept edits localized best, with identity preservation scored by ArcFace. Sole author.
\end{itemize}

\projectEntryPaper{MelanomaNet: Explainable Deep Learning for Skin Lesion Classification}{PyTorch, EfficientNet V2, focal loss, GradCAM++, OpenCV, ISIC 2019}
{https://github.com/suxrobgm/explainable-melanoma}
{https://arxiv.org/abs/2512.09289}
\begin{itemize}
    \item EfficientNet V2 (384$\times$384, focal loss) over all nine ISIC 2019 classes, reaching 0.86 weighted F1 on 25K dermoscopic images.
    \item Splits GradCAM++ attention along the ABCDE criteria. Asymmetry, border irregularity, color variation (K-means), and diameter come straight from the lesion mask, and the paper scores how closely the model's attention agrees with each one. Sole author.
\end{itemize}

%----------PROJECTS----------
\section{Selected Projects}

\projectEntry{LightDepth: Lightweight Monocular Depth Estimation}{PyTorch, ResNet18, U-Net, NYU Depth V2}
{https://github.com/suxrobgm/lightdepth}
\begin{itemize}
    \item ResNet18 encoder and U-Net decoder for single-image depth. It has 42\% fewer parameters than Depth Anything V2 (14.3M vs.\ 24.8M) and runs 72\% faster.
    \item Error on NYU Depth V2 is lower too: AbsRel down 5.3\% and SqRel down 31.1\% against the same baseline.
\end{itemize}

\projectEntry{FSRCNN: Accelerating Super-Resolution Convolutional Neural Network}{PyTorch, mixed-precision training, T91, Set5, Set14, DIV2K}
{https://github.com/suxrobgm/fsrcnn}
\begin{itemize}
    \item Reimplemented FSRCNN (Dong et al., ECCV 2016) at 2$\times$, 3$\times$, and 4$\times$, with upsampling learned end to end.
    \item Recovered the paper's improvement over SRCNN (+1.78\,dB PSNR on Set5, +1.26\,dB on Set14), then ran ablations on the shrinking and mapping layers.
\end{itemize}

\projectEntry{Bookshelf Scanner: Multi-Modal Book Detection and Recognition}{YOLO, Moondream2 VLM, llama.cpp, PyTorch, FastAPI}
{https://github.com/suxrobgm/bookshelf-scanner}
\begin{itemize}
    \item YOLO instance segmentation finds each spine in a shelf photo, and Moondream2 reads the title and author.
\end{itemize}

%----------WORK EXPERIENCE----------
\section{Work Experience}

\experienceEntry{Machine Learning Engineer}{EmTech Care Labs}{Portland, ME}{Jan 2025 -- Present}
% Removed the drafted, unmeasured figures (0.81 AUC, 0.90 F1). Add real numbers back once they are measured.
\begin{itemize}
    \item Built a model that uses daily activity and caregiver logs to flag Alzheimer's residents at risk of decline, so staff can step in earlier.
    \item Wrote an NLP pipeline that turns free-text clinical notes into structured records of symptoms, medications, and care events. Both models run in production on de-identified data under HIPAA.
\end{itemize}

\experienceEntry{Machine Learning Engineer}{AllFactors}{Santa Clara, CA}{Oct 2021 -- Dec 2024}
% Removed the drafted, unmeasured figures (2M listings, 40+ metros, 9% -> 5% valuation error). Add real numbers back once confirmed.
\begin{itemize}
    \item Built the property-valuation and market-forecast models that replaced the company's rules-based system.
    \item Wrote the real-time feature pipeline that fed them (SignalR, SQL Server change tracking) and the Blazor dashboards where users saw the results.
\end{itemize}

\vspace{5pt}\noindent
{\color{lighttext}\textit{Earlier:} Software Engineer, Smart Meal Service, Russia (2020 -- 2021). Game Developer, Pentalight Technology, Malaysia (2020 -- 2021). Freelance (2019 -- 2020). Game Developer, EC Dev Team, Uzbekistan (2016~--~2019).}

%----------TECHNICAL SKILLS----------
\section{Technical Skills}
\begin{description}
    \item[Deep Learning \& CV:] PyTorch, TensorFlow/Keras, OpenCV, CNNs, Vision Transformers, DINOv3, U-Net, EfficientNet, YOLO, contrastive and self-supervised learning, GradCAM++, detection and segmentation, depth estimation
    \item[Mathematics:] Linear algebra, probability and statistics, convex optimization, numerical methods, information theory, hypothesis testing and experimental design
    \item[Machine Learning:] Transfer learning, loss design, regularization, hyperparameter search, cross-validation and error analysis, XGBoost, K-means, NLP
    \item[Research Computing:] Python, NumPy/SciPy, pandas, scikit-learn, Hugging Face, CUDA, mixed-precision and multi-GPU training, Weights \& Biases, Docker, \LaTeX
    \item[Software Engineering:] C\#, C++, TypeScript; .NET, FastAPI, PostgreSQL, Kubernetes, AWS, Azure
\end{description}

\end{document}
